\medskip

On se propose d'étudier la concentration dans le sang d'un médicament ingéré par une personne pour la première fois. Soit $t$ le temps (en heures) écoulé depuis l'ingestion de ce médicament.

On admet que la concentration de ce médicament dans le sang, en gramme par litre de sang, est modélisée par une fonction $f$ de la variable $t$ définie sur l'intervalle $[0~;~+ \infty[$.

\medskip

\textbf{Partie A : lectures graphiques}

\begin{center}
\psset{unit=1.6cm,arrowsize=2pt 3}
\begin{pspicture*}(-0.5,-0.6)(7.5,2.2)
\psgrid[gridlabels=0pt](0,0)(7.5,2.15)
\psaxes[linewidth=1.25pt,labelFontSize=\scriptstyle]{->}(0,0)(0,0)(7.5,2.15)
\psplot[plotpoints=2000,linewidth=1.25pt,linecolor=red]{0}{7.5}{5 x mul 2.71828 x exp div}
\uput[d](6.8,-.2){\footnotesize temps en heures}\uput[r](0,2.1){\footnotesize concentration en g/L}
\end{pspicture*}
\end{center}

On a représenté ci-dessus la courbe représentative de la fonction $f$. Avec la précision permise par le graphique, donner sans justification :

\medskip

\begin{enumerate}
\item Le temps écoulé depuis l'instant de l'ingestion de ce médicament et l'instant où la concentration de médicament dans le sang est maximale selon ce modèle.
\item L'ensemble des solutions de l'inéquation $f(t) \geqslant 1$.
\item La convexité de la fonction $f$ sur l' intervalle [0~;~8].
\end{enumerate}

\medskip

\textbf{Partie B : détermination de la fonction} \boldmath $f$ \unboldmath

\medskip

On considère l'équation différentielle 

\[(E) \::\qquad y' +y = 5\e^{-t},\]

d'inconnue $y$, où $y$ est une fonction définie et dérivable sur l'intervalle $[0~;~+ \infty[$.

On admet que la fonction $f$ est une solution de l'équation différentielle $(E)$.

\medskip

\begin{enumerate}
\item Résoudre l'équation différentielle $(E') : \:\: y' + y = 0$.
\item Soit $u$ la fonction définie sur l’intervalle $[0~;~ +\infty[$ par $u(t) = at\e^{-t}$ avec $a \in \R$.

Déterminer la valeur du réel $a$ telle que la fonction $u$ soit solution de l'équation $(E)$.
\item En déduire l'ensemble des solutions de l'équation différentielle $(E)$.
\item La personne n'ayant pas pris ce médicament auparavant, on admet que $f(0) = 0$.

Déterminer l'expression de la fonction $f$.
\end{enumerate}

\medskip

\textbf{Partie C : étude de la fonction} \boldmath $f$ \unboldmath

\medskip

Dans cette partie, on admet que $f$ est définie sur l'intervalle $[0~;~+\infty[$ par $f(t) = 5t \e^{-t}$.

\medskip

\begin{enumerate}
\item Déterminer la limite de $f$ en $+\infty$.

Interpréter ce résultat dans le contexte de l'exercice.
\item Étudier les variations de $f$ sur l'intervalle $[0~;~ +\infty[$ puis dresser son tableau de variation complet.
\item Démontrer qu’il existe deux réels $t_1$ et $t_2$ tels que $f\left(t_1\right) = f\left(t_2\right) = 1$.

On donnera une valeur approchée à $10^{-2}$ des réels $t_1$ et $t_2$.
\item Pour une concentration du médicament supérieure ou égale à 1 gramme par litre de sang, il y a un risque de somnolence.

Quelle est la durée en heures et minutes du risque de somnolence lors de la prise de ce médicament ?
\end{enumerate}

\medskip

\textbf{Partie D : concentration moyenne}

\medskip

La concentration moyenne du médicament (en gramme par litre de sang) durant la première heure est donnée par :

\[T_m = \displaystyle\int_0^1 f(t) \:\text{d}t\]

où $f$ est la fonction définie sur $[0~;~+\infty[$ par $f(t) = 5t\e^{-t}$.

Calculer cette concentration moyenne.

On donnera la valeur exacte puis une valeur approchée à 0,01 près.

\bigskip
